\documentclass[11pt]{article}

\usepackage[utf8]{inputenc}
\usepackage[T1]{fontenc}
\usepackage{amsmath,amssymb,amsthm,mathtools}
\usepackage[margin=1in]{geometry}
\usepackage{enumitem}
\usepackage{xcolor}
\usepackage[colorlinks=true,linkcolor=blue!60!black,citecolor=blue!60!black,urlcolor=blue!60!black]{hyperref}
\usepackage[nameinlink,capitalize]{cleveref}

\sloppy
\setlength{\emergencystretch}{3em}

\newtheorem{theorem}{Theorem}
\newtheorem{lemma}[theorem]{Lemma}
\newtheorem{corollary}[theorem]{Corollary}
\theoremstyle{definition}
\newtheorem{definition}[theorem]{Definition}
\newtheorem{remark}[theorem]{Remark}

\newcommand{\A}{\mathcal A}
\newcommand{\F}{\mathcal F}
\newcommand{\K}{\mathcal K}
\newcommand{\T}{\mathcal T}
\newcommand{\R}{\mathbb R}
\newcommand{\Q}{\mathcal Q}
\newcommand{\Homp}{\operatorname{Hom}^{+}}
\newcommand{\Ext}{\operatorname{Ext}}
\newcommand{\supp}{\operatorname{supp}}
\newcommand{\brak}[1]{\left\langle #1\right\rangle}
\newcommand{\avg}[2]{\left\llbracket #1 \right\rrbracket_{#2}}

\title{One-Root Planting for the \texorpdfstring{$C_5$}{C5}-Free Class}
\author{}
\date{\today}

\begin{document}
\maketitle

\begin{abstract}
We prove that the class of graphs containing no copy of $C_5$ as a
not-necessarily-induced subgraph is root-plantable at the one-vertex type.
The construction is a local replacement for the independent blow-up used in the
clone-closed case.  Given a rooted $C_5$-free graph $(G,r)$, replace $r$ by a
small positive-density independent cluster and delete all edges inside the old
neighbourhood $N_G(r)$.  The deletion is sparse because $G[N_G(r)]$ is
$P_4$-free, so rooted flag densities change by only $O(\lambda)+O(1/|G|)$; and
the deletion removes exactly the cycles that false-twin cloning of $r$ would
create.  Consequently every one-rooted $C_5$-free limit lies in the closure of
supports of random one-root extensions of unlabelled $C_5$-free limits.
\end{abstract}

\section{Setup}

Let $\T_0$ be the theory of finite simple graphs and let
\[
  \K=\{G:\text{$G$ contains no $C_5$ as a subgraph}\}.
\]
Write $\T_\K$ for the corresponding hereditary constrained theory.  We use the
notation of the root-planting criterion note: for the one-vertex type
$1$, $Q_1$ denotes the space of $C_5$-free one-rooted homomorphisms, and
\[
  S_1=
  \overline{
  \bigcup_{\phi_0\in Q_0}
  \supp(\Ext_1(\phi_0))}
  \subseteq Q_1
\]
is the root-planting set.  Since $\brak{1}=1$, every unlabelled constrained
homomorphism has positive one-vertex type density.

The goal is to prove $S_1=Q_1$.

\section{The local planting construction}

The obstruction to ordinary independent blow-ups is a triangle: blowing up two
vertices of a triangle creates a $5$-cycle.  For one root we can avoid this by
duplicating only the root, and deleting the sparse set of edges inside its
neighbourhood.

\begin{lemma}[Neighbourhood sparsity]\label{lem:nbhd-sparse}
Let $G$ be $C_5$-free and let $r\in V(G)$.  Then $G[N_G(r)]$ contains no copy
of $P_4$ as a subgraph.  In particular,
\[
  e(G[N_G(r)])\le |N_G(r)|\le |V(G)|.
\]
\end{lemma}

\begin{proof}
If $a-b-c-d$ is a path inside $N_G(r)$, then
\[
  r-a-b-c-d-r
\]
is a $5$-cycle in $G$, a contradiction.

It remains only to justify the edge bound.  A connected graph with no $P_4$
subgraph is either a star or a triangle.  Indeed, if it is acyclic, then its
diameter is at most $2$, hence it is a star.  If it contains a cycle, then no
cycle has length at least $4$, and a triangle cannot have any further vertex in
the same connected component, since a shortest path from that vertex to the
triangle would create a $P_4$.  Thus every component has at most as many edges
as vertices, proving the claim.
\end{proof}

\begin{definition}[One-root planting]\label{def:planting}
Let $(G,r)$ be a rooted graph and let $L\ge 1$.  The planted graph
$P_L(G,r)$ is defined as follows.  Remove $r$, add an independent set $R$ of
$L$ new vertices, and keep the old vertices
\[
  U:=V(G)\setminus\{r\}.
\]
Inside $U$, keep every edge of $G-r$ except the edges with both endpoints in
$N_G(r)$.  Finally, join every vertex of $R$ to every vertex of $N_G(r)$, and
to no other vertex of $U$.
\end{definition}

\begin{lemma}[The planted graph is still $C_5$-free]\label{lem:planting-c5-free}
If $G$ is $C_5$-free, then $P_L(G,r)$ is $C_5$-free for every $L\ge1$.
\end{lemma}

\begin{proof}
Put $H=P_L(G,r)$ and let $R$ be the planted root cluster.  Suppose, for a
contradiction, that $H$ contains a $5$-cycle $C$.

If $C$ uses no vertex of $R$, then all its edges are old edges of $G-r$, so it
is a $C_5$ in $G$.

If $C$ uses exactly one vertex $x\in R$, write the cycle as
\[
  x-a-b-c-d-x.
\]
The edges incident with $x$ force $a,d\in N_G(r)$, and the three middle edges
are old edges of $G$.  Hence
\[
  r-a-b-c-d-r
\]
is a $C_5$ in $G$, again impossible.

Finally suppose $C$ uses at least two vertices of $R$.  Since $R$ is
independent and the independence number of $C_5$ is $2$, it uses exactly two,
say $x,y\in R$.  On one of the two arcs between $x$ and $y$ in the cycle there
are two consecutive old vertices, say $b,c$.  Both are adjacent in $H$ to a
vertex of $R$, so both lie in $N_G(r)$; but the edge $bc$ would then be an edge
inside $N_G(r)$, and all such edges were deleted in the construction of $H$.
This contradiction proves the lemma.
\end{proof}

\begin{lemma}[Rooted density approximation]\label{lem:density-approx}
Fix $m\ge1$.  There is a constant $C_m$ such that the following holds.  Let
$(G,r)$ be a rooted $C_5$-free graph on $n$ vertices, let
$H=P_L(G,r)$ with $1\le L\le \lambda n$, and root $H$ at any vertex
$x\in R$.  For every one-rooted flag $F$ with at most $m$ vertices,
\[
  \left|p(F,(H,x))-p(F,(G,r))\right|
  \le C_m\left(\lambda+\frac1n\right).
\]
\end{lemma}

\begin{proof}
Let $\ell=|F|\le m$.  The density $p(F,(H,x))$ is computed by sampling
$\ell-1$ further vertices from $H\setminus\{x\}$.  Couple this with a uniform
sample from $U=V(G)\setminus\{r\}$.

There are two bad events.  First, the sample in $H$ may hit another vertex of
the planted cluster $R$; this has probability at most
$(\ell-1)(L-1)/(n+L-2)=O_m(\lambda+1/n)$.  Second, after the sample has landed
in $U$, it may contain both endpoints of an edge deleted from $G[N_G(r)]$.  By
\Cref{lem:nbhd-sparse}, the number of deleted edges is at most $n$, so this
probability is at most
\[
  \binom{\ell-1}{2}\frac{n}{\binom{n-1}{2}}
  =O_m(1/n).
\]
Outside these bad events, the rooted induced graph seen from $x$ in $H$ is
identical to the rooted induced graph seen from $r$ in $G$: the root has the
same neighbours in $U$, and no edge among the sampled old vertices has been
changed.  Enlarging the implicit constant gives the stated bound.
\end{proof}

\section{Root-plantability at the one-vertex type}

\begin{theorem}\label{thm:c5-one-root}
The $C_5$-free graph class is root-plantable at the one-vertex type:
\[
  S_1=Q_1.
\]
\end{theorem}

\begin{proof}
The inclusion $S_1\subseteq Q_1$ is the general support-passing statement for
hereditary constrained classes.  We prove $Q_1\subseteq S_1$.

Let $\psi\in Q_1$, and let $U$ be a basic neighbourhood of $\psi$, specified by
one-rooted flags $F_1,\ldots,F_s$ of size at most $m$ and a tolerance
$\varepsilon>0$:
\[
  U=
  \{\chi\in Q_1:
    |\chi(F_i)-\psi(F_i)|<\varepsilon\text{ for all }i\}.
\]
Choose a convergent sequence of rooted $C_5$-free graphs
$(G_t,r_t)$ with $|G_t|\to\infty$ and
$p(\,\cdot\,,(G_t,r_t))\to\psi$.

Pick $\lambda>0$ so small that $C_m\lambda<\varepsilon/4$, where
$C_m$ is the constant from \Cref{lem:density-approx}.  For each $t$, put
$L_t=\max(1,\lfloor\lambda |G_t|\rfloor)$ and
\[
  H_t:=P_{L_t}(G_t,r_t).
\]
By \Cref{lem:planting-c5-free}, every $H_t$ is $C_5$-free.  Passing to a
subsequence, assume the unlabelled graphs $H_t$ converge to some
$\phi_0\in Q_0$.

Let $R_t$ be the planted cluster in $H_t$.  The probability that a uniformly
random root of $H_t$ lies in $R_t$ tends to
\[
  \delta=\frac{\lambda}{1+\lambda}>0.
\]
Moreover, for every $x\in R_t$, \Cref{lem:density-approx} and the convergence
of $(G_t,r_t)$ to $\psi$ imply, for all sufficiently large $t$,
\[
  |p(F_i,(H_t,x))-\psi(F_i)|\le\varepsilon/2
  \qquad(1\le i\le s).
\]
Let
\[
  Y_1:=[0,1]^{\F^{1}[\T_0]},
  \qquad
  \widetilde C=
  \{x\in Y_1:
    |x(F_i)-\psi(F_i)|\le\varepsilon/2\text{ for all }i\}.
\]
Then $\widetilde C$ is a closed cylinder in the ambient product space, and
$\widetilde C\cap Q_1\subseteq U$.  The finite random-root distribution of
$H_t$, viewed as a measure on $Y_1$ by recording all rooted flag densities,
assigns mass at least $\delta/2$ to $\widetilde C$ for all large $t$.

By Razborov's random-extension theorem, these finite random-root distributions
converge weakly to $\Ext_1(\phi_0)$ as measures on $Y_1$.  The closed-set form
of the Portmanteau theorem gives
\[
  \Ext_1(\phi_0)(\widetilde C)\ge\delta/2>0.
\]
Since $\Ext_1(\phi_0)$ is concentrated on $Q_1$, its support intersects
$\widetilde C\cap Q_1\subseteq U$.  Since $U$ was an arbitrary neighbourhood
of $\psi$, we have $\psi\in S_1$.
\end{proof}

\begin{remark}
The proof uses two facts special to one root.  First, the only dangerous new
$C_5$s created by cloning the root pass through an old edge inside $N_G(r)$.
Second, those old neighbourhood edges have density $O(1/|G|)$ in rooted
sampling, so deleting them is asymptotically invisible.  For larger types the
analogous construction would have to control edges inside and between several
root-neighbourhood regions; this note does not address that case.
\end{remark}

\end{document}
